\chapter{Parameter Identification} \label{paramIdent}
%
%
This chapter explains how to perform a parameter identification in piezoelectricity using CFS++. 
The \element{analysis} type in your xml-file to choose is \element{paramIdent}, see example below.\\


The parameter identification feature of {\it CFS++} compares simulated and measured data in piezoelectric computations and performs automatically a fit of the data and identifies the complete set (or a subset) of the material parameters involved.
Material characterization can be performed either for a single piezoelectric body or for a compound consisting of a piezoelectric part and additionally one or two adaption (backing) layers which are assumed to be isotropic.

A parameter identification always needs the following steps:
\begin{itemize}
\item Create an initial guess of the material parameters (e.g. parameters from a similar material or IEEE-Standard identification rules). 
\item Select a fitting method, see Table \ref{TableFittingMethods}.
\item Select a measurement quantity, e.g. electrical impedance or mechanical displacements at certain nodes, see Table \ref{TableFittingQuantity}.
\item Determine (carefully) the frequencies at which measurements shall be evaluated (this is actually the most critical part!).
\item Run the identification process. Check for convergence and examine if your results are reasonable.
\end{itemize}
 (All steps are given in details and with examples in Chapter 3 of PhD Thesis {\it Forward and Inverse Problems in Piezoelectricity}, Tom Lahmer, Department of Sensor Technology, Erlangen-Nuremberg, 2008)\\ \newpage

The following analysis section needs to be included in your xml-file:
\begin{verbatim}
<analysis>
  <paramIdent>
    <fittingMethod type="method"/>
    <fittingQuantity type="physical_quantity"/>
    <calcImpedanceCurve> yes/no </calcImpedanceCurve>
    <calcCharge name="surf" neighborRegion="piezo_axi"/>
    <calcMechDisplacementCurve> yes/no </calcMechDisplacementCurve>
    <computeCurveAfterIterationStep> positive_integer </computeCurveAfterIterationStep>
    <numFreq> positive_integer </numFreq>
    <startFreq> positive_integer </startFreq>
    <stopFreq> positive_integer </stopFreq>
    <maxNrIterations> positive_integer </maxNrIterations>
    <stopRes> positive_integer </stopRes>
    <excitationVoltage> positive_integer </excitationVoltage>
    <artDataNoise> positive_integer </artDataNoise>
    <mechDisplAtOppositeNode1> positive_integer </mechDisplAtOppositeNode1>
    <dofOfMechDispl1> positive_integer </dofOfMechDispl1>
  </paramIdent>
</analysis>
\end{verbatim}



\begin{table}
\begin{center}
\caption{Implemented parameter identification methods.}
\label{TableFittingMethods}
\vspace{0.5cm}
\begin{tabular}{lp{0.3\textwidth}p{0.5\textwidth} }
\toprule
{\bf fittingMethodType} & {\bf Method} & {\bf Comments} \\
\midrule
Landweber &  Nonlinear Landweber iteration & robust, however rather slow \\ 
\midrule
minimalError &  modified Landweber iteration with flexible step size choice & robust, usually remarkably faster than Landweber\\ 
\midrule
steepestDescent &  modified Landweber iteration with flexible step size choice  & very similar to the minimal error method \\ 
\midrule
Newton & Newton's method with inexact inner iterations & fast, good convergence properties provided that initial guess is close to solution. Otherwise risk of divergence \\ 
\midrule
NewtonComplex & inexact Newton method. & similar as Newton's method, adapted for simultaneous determination of real and imaginary parts \\
\midrule
leastSquares & simple least-squares error minimizer & converges globally, may become slow after a view steps. It is always a good choice if the initial guess is poor\\
\midrule
confidenceIntervals & computes upper bounds of confidence intervals& provides information about the reliability of the identified parameters (only real parts!). Does not identify anything.\\
\bottomrule
\end{tabular}
\end{center}
\end{table} 

\begin{table}
\begin{center}
\caption{Input quantities.}
\label{TableFittingQuantity}
\vspace{0.5cm}
\begin{tabular}{llp{0.5\textwidth}}
\toprule
 {\bf fittingQuantity} & {\bf Quantity } & {\bf Comments}\\ 
\midrule
logImpedance & $\log(Z)$ & identifies well real and imaginary parts of the parameters \\ \midrule
logAmplitude & $\log(|Z|)$ & identifies well real and imaginary parts of the parameters, is very robust \\ 
\midrule
amplitude & $|Z|$ & less robust, gives very precise results in resonance peaks \\
\midrule
phase & ${\rm arg}(Z)$ & very usefull to identify imaginary parts \\
\midrule
logAmplitudeMech & $\log(|u|)$ & considers logarithm of mechanical displacements \\
\midrule
amplitudeMech & $|u|$ & modulus of mechanical displacement \\
\midrule
displacementMech & $u$ & usually a good choice, very precise \\
\midrule
phaseMech & ${\rm arg}(u)$ & helps to identify imaginary parts \\
\bottomrule
\end{tabular}
\end{center}
\end{table} 

The best way to check the quality of your parameters is the comparison of your simulated impedance curve (or mechanical displacement curve) with your measurements. Setting 
\begin{verbatim}
<calcImpedanceCurve> yes </calcImpedanceCurve>
<calcMechDisplacemetCurve> yes</calcMechDisplacementCurve>
\end{verbatim}
 leads to a calculation of such a curve in the interval $[startFreq,stopFreq]$ with \childelement{numFreq} evaluation points. The output will be a file called {\em imped.dat} ({\em mechDispl.dat}) which can be evaluated with MATLAB/GNUPlot,.. . The file {\em imped.dat} contains the frequency, the computed amplitude, phase, real and imaginary part of the impedance at the surface of the ceramic.  The calculation will be done before and during the parameter identification process (after each  \verb!<computeCurveAfterIterationStep>! iteration step). Comparing the results with the measurements gives hints whether your identification proceeds successfully or not.
Since the computation of the impedance requires the computed surface charges we alway need to provide the tag 
  \verb!<calcCharge name="surf" neighborRegion="piezo_axi"/>! which says on which surface the electric charges are computed. Note, parameter values (except density!) are not taken from the material file {\em e.g. mat.xml}. You have to set them in the file {\em measuredData.dat}, see below.
\\ 

\noindent Further settings: \\
The parameter fitting will be done until some discrepancy principle becomes active, i.e. if the norm of the residual (the misfit between measurement and simulation) falls below the quantity \childelement{$<$stopRes$>$}. If you do not know to which accuracy you can fit the data, select some extremely small value or zero and enforce a stopping of the iterations by setting a maximal number of iteration steps with \childelement{$<$maxNrIterations$>$ 20 - 200 $</$maxNrIterations$>$}.
\\ 
The \childelement{$<$excitationVoltage$>$ 10 $</$excitationVoltage$>$} is supposed to coincide with value of the excitation voltage you specify with your inhomogeneous Dirichlet boundary conditions in the electrostatic part.

Setting artificial data noise makes only sense if you intend to work with synthetically created data (in order to get a feeling how the algorithms behave in case of different amounts of data noise in your measurements). Setting e.g. \childelement{$<$artDataNoise> 0.01 $</$artDataNoise$>$}
would add one per cent random data noise on your impedance or mechanical displacement curve. 

If there are mechanical measurements available (difference of mechanical displacements in one space direction of pairwise facing points) you need to select the appropriate nodes in your mesh and the corresponding degree of freedom (dof).\\
 E.g.  \\ 
\begin{verbatim}
<mechDisplAtNode0> 5 </mechDisplAtNode0>
<mechDisplAtOppositeNode0> 15 </mechDisplAtOppositeNode0>
<dofOfMechDispl0> 0 </dofOfMechDispl0>
\end{verbatim}
 evaluates the difference of the mechanical displacement at the nodes 5 and 15 in $0$ direction (i.e. x- direction). You can select up to three different node pairs.

\subsection{Input Files}
The file  {\em measuredData.dat} provides additional data for the identification process. 
Important: Each line in this file needs to be closed by $'/'$ and all entries need to be separated by $','$.\\ \\
\verb! #  denotes a comment line !\\
\verb! f) 1.0E+06, 2.0E+06, 3.0E+06, ... / ! \\
Numbers after $f)$ are frequencies at which measurements will be read out of the additional input-file {\it mess.dat  ( or messMech.dat)} which contains frequencies, amplitudes and phases (or frequencies, real and imaginary part) of your measurements, e.g. \\
\noindent 
\verb! r) 1.32000E+11, 1.150000E+11, 7.10000E+10, 7.30000E+10,  ... 5.84000E-09,/ ! \\
\verb! i) 0.069500E+11,0.08940E+11, -0.08080E+11,-0.08960E+11,  ... 0.0830E-08,/!\\
\verb! R) 1,  1,  1,  1,  1,  1,   1,  1,  1,   1,/ !\\
\verb! I) 0,  1,  0,  0,  0,  0,  0 ,  1,  0,   1,/ !\\

\noindent The row $r)$ contains initial guesses for the real-valued material parameters in the following order\\ $$c_{11}^E, c_{33}^E, c_{12}^E, c_{13}^E, c_{44}^E, e_{15}, e_{31}, e_{33}, \epsilon_{11}^S, \epsilon_{33}^S.$$ In the next row, after $i)$, the corresponding imaginary parts are listed. The booleans after $R)$ and $I)$ decide which parameters should be updated (identified) during the identification. Sometimes it can be of profit to reduce the set of parameters to those which have a strong impact on the solution in order to save computation time.Depending on your geometry and vibration mode there are always some parameters which are not identifiable at all. In the example above all real parts will be fitted, from the imaginary parts only the three values $c_{33},e_{33}$ and $\epsilon_{33}$. Be aware while determining complex material parameters, that the materialtype is set to complex values in the xml-File, i.e.\\ \noindent \verb!<materialDataType type="imagMaterialParameter"/>! \\
\noindent 


\subsection{Output Files}
Several output files will be created during the run of the program:\\ 
\begin{itemize}
\item \noindent {\em parFinal.dat}, here, you find the computed parameters during the identification. Since they are in the same form as those in the file {\em measuredData.dat} you can simply copy them into this file and reuse them for a new fit, evaluating measurements at different and in particular more frequencies, for example,
\item \noindent {\em parLog.dat}, records the norm of the residual and the development of the parameters,
\item \noindent {\em imped.dat}, impedance curve computed with the current set of parameters, 
\item \noindent {\em mechDispl.dat}, mechanical deflection curve computed with the current set of parameters, 
\item \noindent {\em confInterval.dat} upper bounds of confidence intervals considering the $0.99$ quantile of the $\chi^2$ probability distribution (for details see Chapter 3.3.3 of PhD Thesis, {\it Forward and Inverse Problems in Piezoelectricity}.
\end{itemize}

\subsection{Additional Remarks}
Select \element{$<$gid$/>$} as output quantity. Do not select any store results except the electric charge, otherwise you repeatedly write these values repeatedly the result files which may become extremely large.

Do not forget to set the density in you mat.xml file. This is the only quantity which will not be overwritten by the parameters from the file measuredData.dat

\section{Parameter Identification with Adaption and Backing Layers}
It may be of interest how adaption and backing layers influence the behavior of your transducer and which effective material parameters these layers have. E.g. \\
\begin{verbatim}
<domain geometryType="axi">
  <regionList>
    <region name="piezo_axi" material="PZ27"/>
    <region name="adaption" material="material1"/>
    <region name="backing" material="material2"/>
  </regionList>
</domain>
\end{verbatim}

The fitting feature of {\it CFS++} can handle up to two additional isotropic materials which are characterized only by their mechanical stiffness. The stiffness tensor is composed by the Lam\'{e} parameters $\lambda$ and $\mu$ which can be computed from the elasticity modulus $E$ and the Poisson ratio $\nu$ as follows:
$$\lambda = \frac{\nu E}{(1+\nu)(1-2\nu)},\quad \mu = \frac{E}{2(1+\nu)}.$$
 The extended parameter set is given in the file {\it measuredData.dat} as follows\\
\verb! # piezoelectric constants:! \\
\verb! r) 1.47000E+11, 1.100E+11, 1.05000E+11, 9.37000E+10, ... ,/ !\\
\verb! # first additional material, real parts (mu, lambda):!\\
\verb! a) 7.76e+10, 1.0e+11,/!\\
\verb! # which parameter shall be updated?!\\
\verb! A) 1, 0,/ !\\
\verb! # first additional material, imaginary parts (mu, lambda):!\\
\verb! b) 7.0e+9, 1.5e+9,/!\\
\verb! # which parameter shall be updated?!\\
\verb! B) 1, 0,/ !\\
\verb! #!\\
\verb! # second additional material, real parts (mu, lambda):!\\
\verb! c) 2.6e+10, 5.9e+11,/!\\
\verb! # which parameter shall be updated?!\\
\verb! C) 1, 0,/ !\\
\verb! # second additional material, imaginary parts (mu, lambda):..!\\
\verb! d) 2.2e+9, 2.5e+9,/!\\
\verb! # which parameter shall be updated?! \\
\verb! D) 1, 1,/ !\\
\verb! #!\\

\noindent Create your *.mesh-file in that way that the regions occur in the same order as in the xml-file and that the piezoelectric material is always the first material.



